\section{Formal Interface Excerpts}
\label{app:formal-interface-excerpts}

This appendix records selected Rocq excerpts that determine the meaning of the
main theorem.
They are not intended to replace the full artifact; rather, they expose the
small specification surface that a cryptographic reader should be able to
audit independently.

% Generated by scripts/count-interface-excerpts.py --latex; do not edit counts.
\newcommand{\FormalExcerptLineCount}{724}
\newcommand{\FormalExcerptListingCount}{39}
\newcommand{\FormalExcerptFileCount}{16}
\begin{table}[H]
\centering
\caption{Size of the selected formal-interface excerpts. Counts are physical
Rocq source lines in inclusive \texttt{firstline}--\texttt{lastline} ranges,
including blank and comment lines; typeset wrapping is not counted.}
\label{tab:formal-excerpt-size}
\footnotesize
\begin{tabular}{@{}p{0.78\columnwidth}r@{}}
\toprule
Excerpt category & Source lines \\
\midrule
Scheme interface & 26 \\
Metric charts & 17 \\
Probability-one correctness & 40 \\
Gaussian distribution and KL formula & 20 \\
Product samplers and SSProve bridge & 36 \\
IND-CPA game & 57 \\
IND-CPAD oracles and game & 118 \\
Noise-flooded construction & 45 \\
Concrete IND-CPA reduction & 95 \\
Loss formula & 5 \\
Reduction validity & 6 \\
Miscellaneous specification helpers & 259 \\
\midrule
\textbf{Total} & \textbf{724} \\
\bottomrule
\end{tabular}
\end{table}


The \FormalExcerptListingCount{} listings below select
\FormalExcerptLineCount{} physical source lines from
\FormalExcerptFileCount{} Rocq files;
\Cref{tab:formal-excerpt-size} breaks this count down by category.
The ranges are disjoint, and the count includes blank and comment lines.
This measures the excerpted specification, not the complete human audit
surface. The local specification helpers and top-level theorem statement are
included below; library semantics and assumptions remain on the trust side.
The reduction's running-time audit inspects the displayed code and the costs
of its primitives, rather than adding another source-line category.
The build regenerates the table from the listing ranges with
\texttt{scripts/count-interface-excerpts.py}.

\begingroup
\setmonofont{FreeMono}[Scale=MatchLowercase]

\subsection{Approximate Homomorphic Encryption Scheme Interface}
\label{app:approximate-homomorphic-encryption-interface}

\inputminted[
  fontsize=\scriptsize,
  breaklines,
  linenos,
  firstnumber=25,
  firstline=25,
  lastline=50
]{coq}{rocq-excerpts/theories/Schemes/ApproxFHE.v}

\subsection{Metric Charts}
\label{app:metric-chart-excerpt}

\inputminted[
  fontsize=\scriptsize,
  breaklines,
  linenos,
  firstnumber=105,
  firstline=105,
  lastline=121
]{coq}{rocq-excerpts/theories/Schemes/ApproxFHE.v}

\subsection{Probability-One (Support-Level) Correctness}
\label{app:support-correctness-excerpt}

\inputminted[
  fontsize=\scriptsize,
  breaklines,
  linenos,
  firstnumber=191,
  firstline=191,
  lastline=230
]{coq}{rocq-excerpts/theories/Schemes/ApproxFHE.v}

\subsection{Discrete-Gaussian Distribution and KL Formula}
\label{app:discrete-gaussian-definition-excerpt}

The following excerpts pin down the sampler's mass function.  In particular,
the distribution is not an abstract parameter: it is the normalized
exponential weight below (and is the null subdistribution when the standard
deviation is nonpositive).

\inputminted[
  fontsize=\scriptsize,
  breaklines,
  linenos,
  firstnumber=27,
  firstline=27,
  lastline=29
]{coq}{rocq-excerpts/theories/Probability/DiscreteGaussians/DiscreteGaussian.v}

\inputminted[
  fontsize=\scriptsize,
  breaklines,
  linenos,
  firstnumber=195,
  firstline=195,
  lastline=200
]{coq}{rocq-excerpts/theories/Probability/DiscreteGaussians/DiscreteGaussian.v}

\inputminted[
  fontsize=\scriptsize,
  breaklines,
  linenos,
  firstnumber=226,
  firstline=226,
  lastline=227
]{coq}{rocq-excerpts/theories/Probability/DiscreteGaussians/DiscreteGaussian.v}

\inputminted[
  fontsize=\scriptsize,
  breaklines,
  linenos,
  firstnumber=278,
  firstline=278,
  lastline=282
]{coq}{rocq-excerpts/theories/Probability/DiscreteGaussians/DiscreteGaussian.v}

The equal-variance KL identity consumed by the program logic is itself a
proved theorem, rather than an additional Gaussian assumption.

\inputminted[
  fontsize=\scriptsize,
  breaklines,
  linenos,
  firstnumber=221,
  firstline=221,
  lastline=224
]{coq}{rocq-excerpts/theories/Probability/DiscreteGaussians/DiscreteGaussianKL.v}

\subsection{Product Samplers and SSProve Bridge}
\label{app:discrete-gaussian-product-excerpt}

The construction folds the one-coordinate distribution into an independent
tuple.

\inputminted[
  fontsize=\scriptsize,
  breaklines,
  linenos,
  firstnumber=14,
  firstline=14,
  lastline=26
]{coq}{rocq-excerpts/theories/LibExtras/MathcompExtras/DTuple.v}

\inputminted[
  fontsize=\scriptsize,
  breaklines,
  linenos,
  firstnumber=23,
  firstline=23,
  lastline=24
]{coq}{rocq-excerpts/theories/Constructions/NoiseFlooding.v}

The reduction uses SSProve's integer and vector types.  Its adapter transports
the same one-dimensional distribution into that integer type, then samples
the vector recursively.

\inputminted[
  fontsize=\scriptsize,
  breaklines,
  linenos,
  firstnumber=21,
  firstline=21,
  lastline=22
]{coq}{rocq-excerpts/theories/LibExtras/SSProveExtras/DiscreteGaussian.v}

\begin{samepage}
\inputminted[
  fontsize=\scriptsize,
  breaklines,
  linenos,
  firstnumber=112,
  firstline=112,
  lastline=125
]{coq}{rocq-excerpts/theories/LibExtras/SSProveExtras/DiscreteGaussian.v}
\end{samepage}

The proved bridge identifies that SSProve sampler, after conversion, with the
tuple sampler used by the transformed scheme.

\inputminted[
  fontsize=\scriptsize,
  breaklines,
  linenos,
  firstnumber=205,
  firstline=205,
  lastline=207
]{coq}{rocq-excerpts/theories/Security/IndcpadSimulator.v}

Finally, the addition in the displayed decryption code below is coordinatewise.

\inputminted[
  fontsize=\scriptsize,
  breaklines,
  linenos,
  firstnumber=22,
  firstline=22,
  lastline=23
]{coq}{rocq-excerpts/theories/Schemes/Utils/IntVec.v}

\subsection{IND-CPA Game}
\label{app:ind-cpa-game-excerpt}

\inputminted[
  fontsize=\scriptsize,
  breaklines,
  linenos,
  firstnumber=53,
  firstline=53,
  lastline=109
]{coq}{rocq-excerpts/theories/Schemes/Indcpa.v}

\subsection{IND-CPA with Decryption Oracle and Game}
\label{app:ind-cpa-d-game-excerpt}

\inputminted[
  fontsize=\scriptsize,
  breaklines,
  linenos,
  firstnumber=66,
  firstline=66,
  lastline=130
]{coq}{rocq-excerpts/theories/Schemes/Indcpad.v}

\inputminted[
  fontsize=\scriptsize,
  breaklines,
  linenos,
  firstnumber=138,
  firstline=138,
  lastline=190
]{coq}{rocq-excerpts/theories/Schemes/Indcpad.v}

\subsection{Noise-Flooded Construction}
\label{app:noise-flooded-construction-excerpt}

\inputminted[
  fontsize=\scriptsize,
  breaklines,
  linenos,
  firstnumber=258,
  firstline=258,
  lastline=302
]{coq}{rocq-excerpts/theories/Constructions/NoiseFlooding.v}

\subsection{Concrete IND-CPA Reduction}
\label{app:ind-cpa-reduction-excerpt}

\inputminted[
  fontsize=\scriptsize,
  breaklines,
  linenos,
  firstnumber=224,
  firstline=224,
  lastline=318
]{coq}{rocq-excerpts/theories/Security/IndcpadSimulator.v}

\subsection{Loss Formula}
\label{app:loss-formula-excerpt}

\inputminted[
  fontsize=\scriptsize,
  breaklines,
  linenos,
  firstnumber=39,
  firstline=39,
  lastline=43
]{coq}{rocq-excerpts/theories/Security/NoiseFloodingSecurity/Prelude.v}

\subsection{Reduction Validity}
\label{app:reduction-bound-excerpt}

\inputminted[
  fontsize=\scriptsize,
  breaklines,
  linenos,
  firstnumber=1783,
  firstline=1783,
  lastline=1788
]{coq}{rocq-excerpts/theories/Security/NoiseFloodingSecurity/Final.v}


\subsection{Miscellaneous Specification Helpers}
\label{app:specification-helpers}

The following definitions complete the local dependency chain of the displayed
specification: package interfaces and labels, heap locations, integer-vector
operations and representations, safe table indexing, and the reduction and
security-bound wrappers. The theorem listing repeats the top-level theorem so
that its statement is included in the count. We stop at standard-library,
MathComp, and SSProve primitives, which belong to the trusted semantic base,
and at proved lemmas whose displayed statements are checked by Rocq; we do not
include their proof bodies or proof-only intermediate game definitions.

\inputminted[fontsize=\scriptsize, breaklines, linenos,
  firstnumber=19, firstline=19, lastline=23
]{coq}{rocq-excerpts/theories/Schemes/ApproxFHE.v}

\inputminted[fontsize=\scriptsize, breaklines, linenos,
  firstnumber=52, firstline=52, lastline=98
]{coq}{rocq-excerpts/theories/Schemes/ApproxFHE.v}

\inputminted[fontsize=\scriptsize, breaklines, linenos,
  firstnumber=9, firstline=9, lastline=13
]{coq}{rocq-excerpts/theories/Schemes/Utils/IntVec.v}

\inputminted[fontsize=\scriptsize, breaklines, linenos,
  firstnumber=47, firstline=47, lastline=48
]{coq}{rocq-excerpts/theories/Schemes/Utils/IntVec.v}

\inputminted[fontsize=\scriptsize, breaklines, linenos,
  firstnumber=87, firstline=87, lastline=88
]{coq}{rocq-excerpts/theories/Schemes/Utils/IntVec.v}

\inputminted[fontsize=\scriptsize, breaklines, linenos,
  firstnumber=6, firstline=6, lastline=12
]{coq}{rocq-excerpts/theories/LibExtras/MathcompExtras/ListExtras.v}

\inputminted[fontsize=\scriptsize, breaklines, linenos,
  firstnumber=15, firstline=15, lastline=19
]{coq}{rocq-excerpts/theories/LibExtras/SSProveExtras/ChoiceVector.v}

\inputminted[fontsize=\scriptsize, breaklines, linenos,
  firstnumber=16, firstline=16, lastline=51
]{coq}{rocq-excerpts/theories/Schemes/Indcpa.v}

\inputminted[fontsize=\scriptsize, breaklines, linenos,
  firstnumber=18, firstline=18, lastline=64
]{coq}{rocq-excerpts/theories/Schemes/Indcpad.v}

\inputminted[fontsize=\scriptsize, breaklines, linenos,
  firstnumber=134, firstline=134, lastline=136
]{coq}{rocq-excerpts/theories/Schemes/Indcpad.v}

\inputminted[fontsize=\scriptsize, breaklines, linenos,
  firstnumber=35, firstline=35, lastline=98
]{coq}{rocq-excerpts/theories/Security/IndcpadSimulator.v}

\inputminted[fontsize=\scriptsize, breaklines, linenos,
  firstnumber=320, firstline=320, lastline=321
]{coq}{rocq-excerpts/theories/Security/IndcpadSimulator.v}

\inputminted[fontsize=\scriptsize, breaklines, linenos,
  firstnumber=19, firstline=19, lastline=20
]{coq}{rocq-excerpts/theories/Probability/KL/Core.v}

\inputminted[fontsize=\scriptsize, breaklines, linenos,
  firstnumber=41, firstline=41, lastline=47
]{coq}{rocq-excerpts/theories/Security/NoiseFloodingSecurity/Final.v}

\inputminted[fontsize=\scriptsize, breaklines, linenos,
  firstnumber=522, firstline=522, lastline=533
]{coq}{rocq-excerpts/theories/Security/NoiseFloodingSecurity/Final.v}

\inputminted[fontsize=\scriptsize, breaklines, linenos,
  firstnumber=1953, firstline=1953, lastline=1956
]{coq}{rocq-excerpts/theories/Security/NoiseFloodingSecurity/Final.v}

\inputminted[fontsize=\scriptsize, breaklines, linenos,
  firstnumber=256, firstline=256, lastline=257
]{coq}{rocq-excerpts/theories/Security/NoiseFloodingSecurity/GaussianBasics.v}

\inputminted[fontsize=\scriptsize, breaklines, linenos,
  firstnumber=46, firstline=46, lastline=52
]{coq}{rocq-excerpts/theories/Security/NoiseFloodingSecurity/Prelude.v}

\endgroup
